\section{The theorem and its proof}\label{sec:theorem}

The construction starts from probability measures on finite periodic lattices.
Removing the lattice spacing at fixed physical volume gives the first limit; compactness then gives an infinite-volume limit along a refinement of any prescribed increasing sequence of volumes.
The resulting Euclidean correlations decompose into phases.
One selected phase gives a relativistic theory, while the mixture and every phase admit the same finite operator product expansions.
A logarithmic correction in a degree-four two-point function excludes a specified class of free-field profiles.

We first define the objects and the exact conclusions of the theorem.
The five propositions at the end of this section describe the construction and explain how those conclusions fit together.
The mathematical statements follow the formal source at the fixed revision recorded in~\cite{OpenAIYM}.
Their correspondence with individual declarations is recorded separately; the detailed proofs in Part~II remain to be written.

\subsection{The lattice family}\label{subsec:lattice-family}

Throughout, $\N=\{0,1,\ldots\}$ and the fourth coordinate of $\R^4$ is time.
A gauge group consists of an integer $m\ge1$, a closed subgroup $G\subset U(m)$ and a real linear subspace $\g\subset M_m(\C)$ satisfying
\begin{equation}\label{eq:gauge}
X\in\g\quad\Longleftrightarrow\quad \exp(sX)\in G\text{ for every }s\in\R\,.
\end{equation}
We call $\g$ simple if some $X,Y\in\g$ have $[X,Y]\ne0$ and every real subspace $J\subset\g$ with $[\g,J]\subset J$ equals $0$ or $\g$.
The unitary realization is part of the data; connectedness of $G$ is not required.

For a positive integer $\ell$, let $\Lambda_\ell=(\Z/\ell\Z)^4$.
A link configuration assigns $U(y,\mu)\in G$ to each $y\in\Lambda_\ell$ and $\mu\in\{1,2,3,4\}$.
The plaquette and its Wilson energy are
\begin{equation}\label{eq:wilson-energy}
\begin{aligned}
U_{y;\mu\nu}&=U(y,\mu)U(y+e_\mu,\nu)U(y+e_\nu,\mu)^{-1}U(y,\nu)^{-1},\\
E_\ell(U)&=\sum_{y\in\Lambda_\ell}\sum_{\mu<\nu}\left(1-\frac1m\operatorname{Re}\tr U_{y;\mu\nu}\right)\,.
\end{aligned}
\end{equation}
For normalized Haar measure $h$ on $G$ and $\beta\in\R$, define
\begin{equation}\label{eq:wilson-law}
\mu_{\ell,\beta}(dU)=\frac{e^{-\beta E_\ell(U)}}{Z_{\ell,\beta}}\prod_{y,\mu}h(dU(y,\mu)),
\qquad Z_{\ell,\beta}=\int e^{-\beta E_\ell(U)}\prod_{y,\mu}h(dU(y,\mu))\,.
\end{equation}
Since each plaquette summand lies in $[0,2]$, we have $0\le E_\ell\le12\ell^4$; hence the normalizer is finite and strictly positive.
A gauge transformation $\gamma:\Lambda_\ell\to G$ acts by
\begin{equation*}
U^\gamma(y,\mu)=\gamma(y)U(y,\mu)\gamma(y+e_\mu)^{-1}\,.
\end{equation*}

\begin{definition}[Regular Wilson family]\label{def:family}
A Wilson family $\W$ consists of strictly increasing sequences $L_p,D_N$ of positive integers, real couplings $\beta_{N,p}$, normalized Haar measure $h$, a set $\mathcal T$ of tokens and complex insertions $I_{N,p}(v,y;U)$ for every token $v$, site $y\in\Lambda_{L_pD_N}$ and link configuration $U$ on that lattice.
Put $a_N=D_N^{-1}$ and $\mu_{N,p}=\mu_{L_pD_N,\beta_{N,p}}$.
The family is regular if the following conditions hold.
\begin{enumerate}[label=\textup{Reg\arabic*.}]
\item There exists $b_*>0$ such that, for every $\varepsilon>0$, there exists $N_0$ for which
\begin{equation}\label{eq:running}
\left|\frac{\beta_{N,p}}{\log(1/a_N)}-b_*\right|<\varepsilon
\qquad(N\ge N_0,\ p\in\N)\,.
\end{equation}
\item Every insertion is continuous in $U$.
\item Every insertion is invariant under the gauge transformations above.
\item For each $v\in\mathcal T$ there exists a finite stencil $\Sigma_v\subset\Z^4$ such that, for every $N,p,y$, the insertion depends only on links with both endpoints in $y+(\Sigma_v\bmod L_pD_N)$.
\end{enumerate}
\end{definition}

The fixed stencil shrinks to a point in physical units.
The insertion may depend on both cutoff and volume.
Regularity specifies neither a classical polynomial represented by a token nor a value for the positive slope $b_*$.
The construction will make these choices; they are not additional requirements in this definition.

Write $r_\ell(j)$ for the representative of $j\in\Z/\ell\Z$ in $(-\ell/2,\ell/2]$.
The physical site is $\hat y=a_N(r_{L_pD_N}(y^\mu))_{\mu=1}^4$.
For a token list $v\in\mathcal T^n$ define
\begin{equation}\label{eq:smear}
\begin{aligned}
M_{N,p}(v;y)&=\int\prod_{i=1}^n I_{N,p}(v_i,y_i;U)\,\mu_{N,p}(dU),\\
C_{N,p}(v;f)&=a_N^{4n}\sum_{y\in\Lambda_{L_pD_N}^n} f(\hat y_1,\ldots,\hat y_n)M_{N,p}(v;y)\,.
\end{aligned}
\end{equation}
The separated configuration space is
\begin{equation}\label{eq:conf}
\Conf n=\{x\in(\R^4)^n:x_i\ne x_j\text{ whenever }i\ne j\}\,.
\end{equation}
A separated test $f$ is a globally smooth complex function whose closed support is compact and contained in $\Conf n$.
It is $p$-admissible if its support lies in $((-L_p/2,L_p/2)^4)^n$.
These are joint tests of all points; they need not be products of one-point tests.
Arity zero is included throughout, with empty products equal to $1$.

\begin{definition}[Ordered limit]\label{def:ordered}
For $P:\N\to\N$ and kernels $K_n(v;x)$, the family $\W$ converges along $P$ to $K$ if, for every $n\in\N$, $v\in\mathcal T^n$ and separated test $f$, the test is $P(k)$-admissible eventually in $k$, and there exist numbers $R_p\in\C$ such that
\begin{equation}\label{eq:ordered}
\begin{aligned}
\lim_{N\to\infty}C_{N,p}(v;f)&=R_p
&&\text{($f$ $p$-admissible)},\\
\lim_{k\to\infty}R_{P(k)}&=\int_{(\R^4)^n}f(x)K_n(v;x)\,dx\,.
\end{aligned}
\end{equation}
\end{definition}

The first limit uses the entire ultraviolet sequence at each admissible volume.
The numbers $R_p$ may be assigned arbitrarily at inadmissible volumes.
A strictly increasing $P$ makes each compact test support admissible eventually, since $L_{P(k)}\to\infty$.

\subsection{The continuum objects}\label{subsec:continuum-objects}

Write $x=(\vec x,t(x))$ for spacetime, with Euclidean norm $|x|$ and
\begin{equation}\label{eq:spacetime}
\mathsf m(x)=t(x)^2-|\vec x|^2,
\qquad \overline V_+=\{x:|\vec x|\le t(x)\}\,.
\end{equation}
A vector is spacelike if $\mathsf m(x)<0$.
Let $\mathcal S=\Sch$, with its Schwartz topology.
Hilbert inner products are antilinear in the first variable.

\begin{definition}[Relativistic data]\label{def:relativistic}
A relativistic datum consists of a complex Hilbert space $\mathcal H$, a complex vector space $\mathcal L$ of labels, a linear subspace $\mathcal D\subset\mathcal H$, a vector $\Omega\in\mathcal D$ and fields $\Phi(\ell,f):\mathcal D\to\mathcal D$ linear over $\C$ in $\ell\in\mathcal L$ and $f\in\mathcal S$.
It includes an antilinear map $\ell\mapsto\ell^*$, a strongly continuous unitary translation representation $T(b)$ and a projection-valued Borel measure $E$ on $\R^4$.
Its diagonal finite measures satisfy $\mu_\psi(A)=\langle\psi,E(A)\psi\rangle$ on Borel sets.

There is also a strongly continuous unitary representation $U$ of a topological group $\mathcal P$, identified by a homeomorphism $\kappa$ with the proper orthochronous Poincar\'e group.
Writing $\kappa(g)=(\Lambda_g,b_g)$, this identification respects identity and the multiplication
\begin{equation}\label{eq:poincare}
(\Lambda,b)(\Lambda',b')=(\Lambda\Lambda',\Lambda b'+b)\,.
\end{equation}
Here $\Lambda$ is invertible, has determinant $1$, preserves $\mathsf m$ and maps $\overline V_+$ onto itself.
Pure translations $\tau(b)=\kappa^{-1}(1,b)$ satisfy $U(\tau(b))=T(b)$.
For each $g\in\mathcal P$ there is a complex linear bijection $\rho(g)$ of $\mathcal L$.
The action on tests is $(g\cdot f)(x)=f(\Lambda_g^{-1}(x-b_g))$.
\end{definition}

A projection-valued Borel measure here is normalized by $E(\R^4)=1$, has $E(\varnothing)=0$, satisfies $E(A)E(B)=E(A\cap B)$ and $E(A)^*=E(A)$ on Borel sets, and is countably additive in the strong operator topology.
One can assign arbitrary operators to non-Borel sets, which none of the assertions uses.
The label map ${}^*$ is not required to be an involution.
No continuity or finite-dimensionality condition on $\rho$ is included in the data.

\begin{definition}[Wightman conditions]\label{def:wightman}
The relativistic datum satisfies the following conditions, with $p\cdot b$ the Euclidean pairing.
\begin{enumerate}[label=\textup{W\arabic*.},start=0]
\item For every $\psi\in\mathcal H$ and $b\in\R^4$,
\begin{equation}\label{eq:spectrum}
\langle\psi,T(b)\psi\rangle=\int e^{ip\cdot b}\,\mu_\psi(dp),
\qquad E(\R^4\setminus\overline V_+)=0\,.
\end{equation}
Every vector fixed by all $T(b)$ belongs to $\C\Omega$.
\item The vacuum has norm one, $\mathcal D$ is dense, and
\begin{equation}\label{eq:domain}
\mathcal D=\Span_{\C}\{\Phi(\ell_1,f_1)\cdots\Phi(\ell_j,f_j)\Omega:j\in\N,\ \ell_i\in\mathcal L,\ f_i\in\mathcal S\}\,.
\end{equation}
For every $\ell\in\mathcal L$ and $\psi\in\mathcal D$, the map $f\mapsto\Phi(\ell,f)\psi$ is continuous from $\mathcal S$ into $\mathcal H$ in norm.
For every $\ell,f$ and $\psi,\chi\in\mathcal D$,
\begin{equation}\label{eq:adjoint}
\langle\Phi(\ell,f)\psi,\chi\rangle=\langle\psi,\Phi(\ell^*,\bar f)\chi\rangle\,.
\end{equation}
\item The maps $\rho$ form a representation: $\rho(1)=1$ and $\rho(gg')=\rho(g)\rho(g')$.
For every $g\in\mathcal P$, $U(g)\mathcal D\subset\mathcal D$ and $U(g)\Omega=\Omega$.
For every $g,\ell,f$ and $\psi\in\mathcal D$,
\begin{equation}\label{eq:covariance}
U(g)\Phi(\ell,f)\psi=\Phi(\rho(g)\ell,g\cdot f)U(g)\psi\,.
\end{equation}
\item If $\mathsf m(x-y)<0$ for every $x\in\supp f$ and $y\in\supp f'$, then, for every $\ell,\ell'\in\mathcal L$ and $\psi\in\mathcal D$,
\begin{equation}\label{eq:locality}
\Phi(\ell,f)\Phi(\ell',f')\psi=\Phi(\ell',f')\Phi(\ell,f)\psi\,.
\end{equation}
\end{enumerate}
\end{definition}

The word span is the domain itself, and includes the vacuum as the empty word.
The symmetry group is the Poincar\'e group itself rather than its double cover.
The sign convention in~\eqref{eq:spectrum} is converted to the usual Minkowski pairing by reversing spatial momentum; the forward cone is unchanged.

\begin{definition}[Hierarchy]\label{def:hierarchy}
A hierarchy over $\W$ contains the following data.
\begin{enumerate}[label=\textup{H\arabic*.}]
\item A relativistic datum of Definition~\ref{def:relativistic}.
\item For every $d\in\N$, a finite-dimensional complex normed space $F_d$ and a complex linear map $\iota_d:F_d\to\mathcal L$.
Set $\mathbb F=\bigsqcup_{d\in\N}F_d$, $\deg(d,A)=d$ and $\iota(d,A)=\iota_d A$.
\item A projection $\pi:\mathcal T\to\mathbb F$.
\item A compact metric space $X$ with its Borel measure $\nu$ and a selected point $\xi_*\in X$.
\item For every $n\in\N$ and $s\in X\sqcup\{\mix\}$, a correlation function $S^s(a;x)\in\C$ on $\mathbb F^n\times(\R^4)^n$.
\item For every $n\in\N$, a continuation $\mathsf F_n(a;z,\vec x)\in\C$ on $\mathbb F^n\times\C^n\times(\R^3)^n$.
\item For every $j\in\N$, a finite channel set $\Ch j$ and fields $O_{j,\kappa}\in F_j$.
For every finite index set $I$, $a\in\mathbb F^I$ and $\kappa\in\Ch j$, there is a coefficient function $c_{j,\kappa}(a;x,o)$ on $(\R^4)^I\times\R^4$.
\item An element $\Theta\in F_4$.
\end{enumerate}
\end{definition}

The degree is an index used in the product estimates below.
The definition imposes no inclusions between different $F_d$, no prescribed curvature interpretation and no compatibility of the $\iota_d$ across degrees.
The functions $S^s$ are defined even on collision diagonals, but the conclusions constrain them only at separated points.

\subsection{Reconstruction and phases}\label{subsec:reconstruction}

The mixture will be the limit of lattice moments.
The point $\xi_*$ instead selects the Euclidean correlations that continue to the relativistic fields.
These two roles are related by a phase average, not by identifying the two states.

\begin{definition}[Reconstruction properties]\label{def:reconstructed}
A hierarchy has the reconstruction properties if the following hold.
\begin{enumerate}[label=\textup{R\arabic*.}]
\item Its relativistic datum satisfies W0--W3.
\item For every nonzero $\ell\in\mathcal L$, there exists $f\in\mathcal S$ with $\Phi(\ell,f)\ne0$.
\item Each $\iota_d$ is injective.
\item Some $a\in\mathbb F$ and $f\in\mathcal S$ satisfy $\Phi(\iota(a),f)\ne0$.
\item The map $a\mapsto\iota(a)$ from $\mathbb F$ onto $\mathcal L$ is surjective.
\item The map $\pi:\mathcal T\to\mathbb F$ is surjective.
\item The measure $\nu$ has total mass one.
\item Every nonempty open subset of $X$ has positive $\nu$-measure.
\item For every state $s$, arity $n$, list $a\in\mathbb F^n$, $x\in\Conf n$ and permutation $\sigma$,
\begin{equation}\label{eq:permutation}
S^s((a_{\sigma(i)})_i;(x_{\sigma(i)})_i)=S^s(a;x)\,.
\end{equation}
\item For every $s,n,a$, the function $S^s(a;\cdot)$ is smooth on $\Conf n$.
\item For every $n,a$ and $x\in\Conf n$, the function $\xi\mapsto S^\xi(a;x)$ is continuous.
\item For every $n,a$ and $x\in\Conf n$,
\begin{equation}\label{eq:mixture}
\int_X S^\xi(a;x)\,\nu(d\xi)=S^\mix(a;x)\,.
\end{equation}
\item For every $n,a$ and fixed real spatial configuration $\vec x$, the function $z\mapsto\mathsf F_n(a;z,\vec x)$ is holomorphic on
\begin{equation}\label{eq:tube}
\mathcal Z_n=\{z\in\C^n:\operatorname{Im}z_1<\cdots<\operatorname{Im}z_n\}\,.
\end{equation}
\item For every $n,a$ and real configuration $x$ with $t(x_1)<\cdots<t(x_n)$,
\begin{equation}\label{eq:euclidean-root}
\mathsf F_n(a;(it(x_i))_i,(\vec x_i)_i)=S^{\xi_*}(a;x)\,.
\end{equation}
\item For every $n,a$, real heights $h_1<\cdots<h_n$ and tests $f_1,\ldots,f_n\in\mathcal S$,
\begin{multline}\label{eq:boundary}
\lim_{r\downarrow0}\int_{(\R^4)^n}\mathsf F_n(a;(t(x_i)+irh_i)_i,(\vec x_i)_i)\prod_{i=1}^n f_i(x_i)\,dx\\
=\langle\Omega,\Phi(\iota(a_1),f_1)\cdots\Phi(\iota(a_n),f_n)\Omega\rangle\,.
\end{multline}
\end{enumerate}
\end{definition}

In~\eqref{eq:boundary} we retain the formal source's convention that the integral of a nonintegrable function is zero.
No separate positive-height integrability assertion is part of R15.
The integrals in the lattice limit and~\eqref{eq:mixture} are ordinary absolutely convergent integrals, by compact support and smoothness in the first case and by phase continuity and compactness in the second.

\subsection{Products and the trace test}\label{subsec:products}

The product expansion replaces a cluster by finitely many single fields at its center.
One choice of coefficients must work with every list of distant spectators and in every state.
The threshold on the cluster radius, however, may depend on the fixed spectator fields.

For finite index sets $I,J$, write $S^s(a,x;b,y)$ for the correlation on the concatenated lists, enumerated in one fixed order on $I\sqcup J$.
For $r>0$, $M\ge0$ and $\eta>0$, call $(x,o)$ an $(r,M,\eta)$-cluster if
\begin{equation}\label{eq:cluster}
|x_i-o|\le Mr\quad(i\in I),
\qquad |x_i-x_{i'}|\ge\eta r\quad(i\ne i')\,.
\end{equation}
For $c>0$, call $y$ $c$-clear of $o$ if
\begin{equation}\label{eq:clear}
|o-y_j|\ge c\quad(j\in J),
\qquad |y_j-y_{j'}|\ge c\quad(j\ne j')\,.
\end{equation}
Set $w(a)=\sum_{i\in I}\deg a_i$.

\begin{definition}[Operator product expansion]\label{def:ope}
The hierarchy has the OPE property if both assertions below hold.
\begin{enumerate}[label=\textup{O\arabic*.}]
\item For every finite $I$, $a\in\mathbb F^I$, $M\ge0$, $\eta,\varepsilon>0$, $j\in\N$ and $\kappa\in\Ch j$, there exists $\delta>0$ such that, for every $0<r<\delta$ and every $(r,M,\eta)$-cluster,
\begin{equation}\label{eq:ope-coefficient}
|c_{j,\kappa}(a;x,o)|\le r^{j-w(a)-\varepsilon}\,.
\end{equation}
\item For every finite $I,J$, $a\in\mathbb F^I$, $b\in\mathbb F^J$, $M\ge0$, $\eta,c,\varepsilon>0$ and $D\in\N$, there exists $\delta>0$ such that, for every $0<r<\delta$, every $(r,M,\eta)$-cluster, every $c$-clear spectator configuration and every $s\in X\sqcup\{\mix\}$,
\begin{equation}\label{eq:ope-remainder}
\left|S^s(a,x;b,y)-\sum_{j=0}^D\sum_{\kappa\in\Ch j}
 c_{j,\kappa}(a;x,o)S^s((j,O_{j,\kappa}),o;b,y)\right|\le r^{D+1-w(a)-\varepsilon}\,.
\end{equation}
\end{enumerate}
\end{definition}

The coefficient functions are those already chosen in H7.
They do not depend on spectators, state or truncation degree.
Both estimates have constant $1$, with the dependence absorbed into $\delta$.
For small enough $r$, $Mr<c$ ensures that all points in the tested correlations are separated.
These are finite-truncation estimates; they assert neither an infinite convergent expansion nor uniqueness of its coefficients.

The last conclusion concerns a single two-point function in the selected phase:
\begin{equation}\label{eq:trace-pair}
G_\Theta(r)=S^{\xi_*}((4,\Theta),(4,\Theta);-re_4/2,re_4/2)\,.
\end{equation}
To specify the comparison class, a mode consists of $\mathsf m\ge0$, integers $j,q\in\N$ and a continuous real function $h$ on $\R^3$ satisfying $h(\lambda p)=\lambda^jh(p)$ for every $\lambda>0$.
Set
\begin{equation}\label{eq:mode}
E_{\mathsf m}(p)=\sqrt{\mathsf m^2+|p|^2},
\qquad K_{\mathsf m,q,h}(r)=\int_{\R^3}\frac{E_{\mathsf m}(p)^qh(p)}{2E_{\mathsf m}(p)}e^{-rE_{\mathsf m}(p)}\,dp\,.
\end{equation}
The integrand can be assigned value zero at a zero denominator, a set of measure zero.
A finite Wick profile is
\begin{equation}\label{eq:wick}
W(r)=\sum_{\alpha=1}^A c_\alpha\prod_{\beta=1}^{n_\alpha}
 K_{\mathsf m_{\alpha\beta},q_{\alpha\beta},h_{\alpha\beta}}(r),
\qquad c_\alpha\in\C,\quad A,n_\alpha\in\N\,.
\end{equation}
Empty products are allowed.

\begin{definition}[Trace test]\label{def:trace}
The hierarchy passes the trace test if the following hold.
\begin{enumerate}[label=\textup{T\arabic*.}]
\item For every $f\in\mathcal S$, $\langle\Omega,\Phi(\iota_4\Theta,f)\Omega\rangle=0$.
\item The number $A_G=3\dim_\R\g/\pi^4$ is positive.
\item The following limit holds in $\C$ over all positive radii:
\begin{equation}\label{eq:trace-law}
\lim_{r\downarrow0}r^8\bigl(\log(1/r)\bigr)^2G_\Theta(r)=\frac{3\dim_\R\g}{\pi^4}\,.
\end{equation}
\item For every finite Wick profile $W$ and every $\delta>0$, there exists $r\in(0,\delta)$ with $G_\Theta(r)\ne W(r)$.
\end{enumerate}
\end{definition}

The public trace test concerns the chosen degree-four field and the time axis.
Its definition alone supplies no stress tensor.
The construction identifies this field with a centered stress trace, as used in the proof of Proposition~\ref{prop:trace} below.

\subsection{The main theorem}\label{subsec:main-theorem}

A hierarchy over $\W$ has the main properties along $P$ if:
\begin{enumerate}[label=\textup{P\arabic*.}]
\item $\W$ converges along $P$ in Definition~\ref{def:ordered} to $K_n(v;x)=S^\mix(\pi(v_1),\ldots,\pi(v_n);x)$;
\item it has properties R1--R15;
\item it has properties O1--O2;
\item it has properties T1--T4.
\end{enumerate}
These four assertions refer to one hierarchy, including the very same spectator fields in P1 and in the OPE.

\begin{theorem}\label{thm:main}
For every gauge group $(m,G,\g)$ with $\g$ simple, there exists a regular Wilson family $\W$ such that, for every strictly increasing $p:\N\to\N$, there exist a strictly increasing $q:\N\to\N$ and a hierarchy over $\W$ with the main properties along $p\circ q$.
\end{theorem}

The lattice data are chosen before $p$.
The refinement, field spaces, phases, selected relativistic theory, product coefficients and trace field may depend on $p$.
The same refinement works for all arities, tokens and separated tests.
The lattice limit is the mixture; the relativistic theory is attached to $\xi_*$.

The theorem does not assert uniqueness of the infinite-volume limit, universality with respect to lattice data, or convergence along simultaneous cutoff and volume limits.
It does not assert a mass gap, clustering of the mixture, agreement of phases, or that an individual phase is a lattice-state limit.
It gives no contact extension across collisions, and excludes only the finite profiles in~\eqref{eq:wick}, not all composites of Gaussian fields or continuous-mass generalized free fields.
No Euclidean covariance or reflection-positivity assertion is included among P1--P4; those properties are established internally in the construction.

\subsection{The principal propositions}\label{subsec:principal}

The following propositions assert one joint choice of objects.
For each gauge group there exists one family satisfying Proposition~\ref{prop:uv}; for each increasing $p$ there exist one increasing $q$ and one hierarchy satisfying Propositions~\ref{prop:limit}--\ref{prop:trace} together.
The latter assertions concern this chosen family and hierarchy.
They are not assertions about an arbitrary regular family or about an arbitrary table of smooth correlations.
This joint choice is essential: separate existence of objects satisfying the separate conclusions would not imply Theorem~\ref{thm:main}.

\paragraph{Removing the ultraviolet cutoff.}
The short-distance integration must control the observables as well as the density.
At fixed volume the bare coupling changes with the cutoff, so comparison at a fixed bare coupling cannot establish convergence.
Instead, each cutoff trajectory is tuned to one common effective coupling at unit scale.

\begin{proposition}[Ultraviolet construction]\label{prop:uv}
For every gauge group with simple $\g$, there exists a regular Wilson family $\W$ in the joint construction with the following property.
For every $p,n\in\N$ and $v\in\mathcal T^n$, there exists a function $K_{p,v}$ smooth on
\begin{equation}\label{eq:periodic-separated}
\{x\in(\R^4)^n:x_i-x_j\notin L_p\Z^4\text{ whenever }i\ne j\}
\end{equation}
such that, for every $p$-admissible separated test $f$,
\begin{equation}\label{eq:fixed-volume}
\lim_{N\to\infty}C_{N,p}(v;f)=\int f(x)K_{p,v}(x)\,dx\,.
\end{equation}
\end{proposition}

\emph{Idea of the proof.}
An exactly normalized, gauge-covariant block kernel preserves the partition integral.
Gaussian integration around a stationary background and a connected polymer expansion estimate the small-field part, while an energy lower bound controls the retained large-field part.
The marginal curvature-square coefficient is absorbed into the running coupling by an exact identity of densities.
The choices of block ratios and entry couplings make the accumulated coupling error smaller than the total logarithmic scale.
Marked insertions evolve through these same steps, with a fixed reference normalization for each finite degree.
Their uniform tails allow the fixed-level limit to imply the full fixed-volume limit: first make the tail small, then compare large cutoffs at that fixed level.

\gap{The Part II proof must supply the nonlinear-step estimates, attach the one-loop coefficient to the actual stationary determinant, and prove convergence of the tuned fixed-level data with uniform observable tails.
The stated existence and full-spacing limit have been checked at the formal assembly; these analytic proofs have not yet been written here.\sourcefootnote{\nolinkurl{Construction.lean}, lines 15--40; \nolinkurl{HierarchyBridge.lean}, lines 16--26; \nolinkurl{YM/FullVolume/Profiles.lean}, lines 25--47.}}

\paragraph{Taking infinite volume and resolving the mixture.}
Local estimates independent of the torus permit a common compactness extraction.
Their reflection-positive Hilbert space may have more than one ground state.
Compressing bounded slab insertions to that ground space gives the algebra whose characters describe the phases.

\begin{proposition}[Continuum limit and phases]\label{prop:limit}
For the jointly chosen family of Proposition~\ref{prop:uv} and every strictly increasing $p:\N\to\N$, there exist a strictly increasing $q:\N\to\N$ and hierarchy data H1--H8 such that P1 and R7--R12 hold.
The same hierarchy satisfies Propositions~\ref{prop:reconstruction}--\ref{prop:trace}.
\end{proposition}

\emph{Idea of the proof.}
Volume-uniform chessboard bounds control separated correlations.
A positive transfer kernel and complex interpolation give equicontinuity, and countable extraction gives a common refinement for all finite label bases and compact sets.
Lattice summation by parts and corrected higher-degree representatives identify positional derivatives.
The full fixed-volume limits of Proposition~\ref{prop:uv} identify the inner limits, removing the auxiliary ultraviolet refinement.

For the phase decomposition, powers of a slab operator are bounded with an exterior-word factor incurred once, independently of the number of powers.
Taking roots removes this factor and gives a bounded operator.
Continuous rotations, proved by stress-shell identities and comparison with free Maxwell fields, allow the time and a spatial axis to be exchanged.
This exchange proves commutativity of the ground-space algebra.
Its compact spectrum and the measure representing the vacuum give~\eqref{eq:mixture}.

\gap{The first limit section must state the actual ultraviolet estimates needed for compactness and for corrected derivative representatives.
The later symmetry and phase sections must prove the exterior-flux estimate and commutativity before using the character decomposition; compactness alone does not give these conclusions.\sourcefootnote{\nolinkurl{YM/Extraction/Limits.lean}, lines 26--43; \nolinkurl{HierarchyBridge.lean}, lines 44--57 and 82--119; the source phase construction is \nolinkurl{YM/Phase/Mixture.lean}.}}

\paragraph{Continuing a selected phase.}
A phase with one-dimensional ground space is selected without imposing an extra condition on its trace correlations.
The continuation is built from bounded slab insertions separated by positive heat evolution.
Bounds as the slabs become thin control its boundary values.

\begin{proposition}[Relativistic reconstruction]\label{prop:reconstruction}
For the jointly chosen hierarchy in Proposition~\ref{prop:limit}, properties R1--R6 and R13--R15 hold, with the same $F_d$, $\pi$, $\iota_d$, $\xi_*$ and $S^{\xi_*}$.
\end{proposition}

\emph{Idea of the proof.}
Polynomial thin-slab bounds yield boundary distributions.
Positivity realizes their field words on a Hilbert space; prepending a field defines the operators on the word span.
Continued rotation Ward identities give Lorentz covariance, and the positive-energy half-space together with that covariance gives forward-cone spectrum.
Agreement of the two analytic orderings at a spacelike adjacent pair proves local commutativity.
Faithfulness follows by recovering the Euclidean correlations from a vanishing continued field.
The quotient by all separated null correlations makes each $\iota_d$ injective, and compatible finite-degree fields cover the label space.
The unit field and faithfulness supply a nonzero field action.

\gap{The reconstruction section must prove the joint boundary estimates, the spacelike analytic gluing and their application to the chosen phase.
The public properties, including norm continuity of smearing and the precise product-test boundary limit, have been checked in the formal conversion; the analytic arguments are still deferred.\sourcefootnote{\nolinkurl{PublicTheory.lean}, lines 93--319; \nolinkurl{HierarchyBridge.lean}, lines 82--133; boundary and locality proofs are organized in \nolinkurl{YM/Work/C01.lean} through \nolinkurl{C29.lean}.}}

\paragraph{Expanding a local cluster.}
The lattice evolution expresses a cluster as a local insertion plus a remainder.
A finite family of separating correlations reads its coefficients degree by degree.
An operator-norm remainder estimate in the commutative ground algebra then applies to every character and the mixed vacuum at once.

\begin{proposition}[Common operator product expansion]\label{prop:ope}
The fields and coefficient functions H7 of the jointly chosen hierarchy satisfy O1--O2 of Definition~\ref{def:ope}.
One coefficient system serves every spectator list, state and finite truncation degree.
\end{proposition}

\emph{Idea of the proof.}
Integrate the clustered product from its scale toward the spectator scale, retaining its components of degree at most $D$.
The remainder gains the next power of the cluster radius, with an arbitrarily small exponent loss absorbing the logarithmic factors.
Choose the separating correlations once, before any spectator is prescribed.
Apply the resulting algebra norm bound before evaluating in a phase; this gives the uniformity in $s$ in~\eqref{eq:ope-remainder}.
The permitted dependence of the radius threshold on the spectator fields is retained.

\gap{The product section must prove the marked cluster estimate and the passage from the chosen separating correlations to the algebra-norm remainder.
The exact eventual-radius quantifiers of O1--O2 have been checked against the formal statements, but these estimates have not yet been reproduced.\sourcefootnote{\nolinkurl{ChallengeSpec.lean}, lines 58--80; \nolinkurl{HierarchyBridge.lean}, lines 68--80; \nolinkurl{YM/OPE/Orbits/Kernel.lean}, lines 62--77.}}

\paragraph{Detecting the logarithmic correction.}
The construction supplies a conserved stress tensor and centers its trace in the selected phase.
Short-distance comparison with free Maxwell fields determines the trace's leading behavior along suitable scales.
Positivity of its spectral measure and the ability to fit a subsequence of every sequence of radii give the full right limit.

\begin{proposition}[Trace law and finite-profile exclusion]\label{prop:trace}
The field $\Theta\in F_4$ of the jointly chosen hierarchy satisfies T1--T4 of Definition~\ref{def:trace} in the selected theory of Proposition~\ref{prop:reconstruction}.
\end{proposition}

\emph{Idea of the proof.}
The trace comparison contributes one inverse running coupling at each insertion.
Its logarithmic running therefore produces the inverse square logarithm in~\eqref{eq:trace-law}.
The one-loop coefficient enters both the trace normalization and the coupling; their ratio cancels from the final amplitude.
Thus that amplitude alone does not check the numerical beta-function coefficient.

The positive spectral Laplace representation makes $G_\Theta$ nonincreasing.
Comparison at fitted separations on either side of a prescribed radius then gives the same asymptotic there.
Fitting a subsequence of any proposed counterexample sequence proves the full-radius limit.
Finite products and sums of the kernels in~\eqref{eq:mode} instead have expansions with integer powers and nonnegative powers of the logarithm.
The finite-order asymptotic estimates for those expansions preclude a finite nonzero limit after multiplication by $r^8\log^2(1/r)$, which proves T4.

\gap{The trace section must prove the stress normalization, flexible-scale comparison and spectral sandwich, and the finite mass-shell expansion used in the exclusion.
The trace amplitude and the conversion of the comparison class have been checked directly; no statement about continuous-mass generalized free fields is inferred.\sourcefootnote{\nolinkurl{TraceBridge.lean}, lines 14--44; \nolinkurl{YM/TraceExtension/Verse/Cape.lean}, lines 53--88; \nolinkurl{YM/TraceExtension/Wick/Log.lean}, lines 96--108.}}

\begin{proof}[Deduction of Theorem~\ref{thm:main} from the principal propositions]
Fix the gauge group and choose the family of Proposition~\ref{prop:uv}, including its insertions and positive regularity slope.
For a prescribed strictly increasing $p$, take the common $q$ and hierarchy of Proposition~\ref{prop:limit}.
This gives P1 and R7--R12 on that refinement.
Proposition~\ref{prop:reconstruction} supplies the remaining reconstruction properties, hence P2.
Propositions~\ref{prop:ope} and~\ref{prop:trace} give P3 and P4 for those same objects.
\end{proof}

The joint existence required here follows in the formal proof by choosing the block-kernel penalty before the family, specializing it to $1$, and using the full ultraviolet sequence with the prescribed volume sequence.
The detailed construction then puts the limit table, selected theory, phases, product coefficients and trace on one refinement.
This checks the logical reduction and the identities of its objects; it does not replace the deferred analytic proofs above.\sourcefootnote{\nolinkurl{Solution.lean}, lines 4--10, specializes \texttt{YangMills.constructed}; the common witnesses and all four properties are assembled in \nolinkurl{Construction.lean}, lines 15--40.}
